\paragraph{Statement A.5.5 (Sharp-arc mean square; complete candidate proof under review).}
For every $\epsilon>0$,
\begin{equation}\label{A:sharp-mean-square}
 \int_X^{2X}|R(x)-C_0x|^2\,dx\ll_\epsilon X^{2+\epsilon}.
\end{equation}
\begin{proof}
We extend the uniform transform of Statement A.4.4 to a negative
Mellin line and then apply the geodesic large sieve to its actual
arc coefficients. This checks the additional step needed to use that
large sieve here.

\emph{1. The geometric correction on a negative line.}
In \eqref{A:F-spherical-integral} put
$p=(s+i\kappa)/2$, $B(x)=\cosh(2\delta)+x\sinh(2\delta)$,
and $f(x)=B(x)^{-p}$, where $\delta=\operatorname{arsinh}(1/A)$.
For $\operatorname{Re}s>0$ the identity
\[
 (1-x^2)^{s/2-1}
 =\frac{s+1}{s}(1-x^2)^{s/2}
       -\frac1s\frac d{dx}\{x(1-x^2)^{s/2}\}
\]
and integration by parts give
\begin{equation}\label{A:F-regularized}
 F_{s,\kappa}(A)=
 \frac{\Gamma((s+3)/2)}{\sqrt\pi\Gamma((s+2)/2)}
 \int_{-1}^1(1-x^2)^{s/2}
       \left(f(x)+\frac{x f'(x)}{s+1}\right)dx.
\end{equation}
The right side converges absolutely and is analytic for
$\operatorname{Re}s>-2$, apart from its displayed parameter
singularities. Therefore it continues the left side to the range
$-1<\operatorname{Re}s<0$ used below.

Fix $-1<\sigma<0$, write $s=\sigma+i\tau$,
$U=1+|\tau|$ and $V=U+|\kappa|$. The gamma prefactor is
$O_\sigma(U^{1/2})$, $|s+1|^{-1}\ll_\sigma U^{-1}$,
and $|f'(x)|\ll_\sigma V/A$. Hence
\begin{equation}\label{A:F-negative-large}
 |F_{s,\kappa}(A)-1|
       \ll_\sigma U^{1/2}\left(1+\frac{V}{UA}\right).
\end{equation}
If $A\geq V$, differentiate the integrand in
\eqref{A:F-regularized} twice with respect to $\delta$ on
$0\leq\delta\leq\operatorname{arsinh}(1/A)$.
Its first derivative at zero is odd in $x$, so it integrates to
zero. Its second derivative is $O_\sigma(V^2)$ times the
integrable absolute weight: the possible third power of $p$
from $f'$ is multiplied by $\delta/(s+1)$ and is bounded by
$O_\sigma(V^2)$ because $V\delta\leq1$.
Taylor's integral remainder therefore gives
\begin{equation}\label{A:F-negative-small}
 |F_{s,\kappa}(A)-1|\ll_\sigma U^{1/2}V^2/A^2
 \qquad(A\geq V).
\end{equation}
Using the count $\#\{A_z\leq y\}\ll y$ with its ideal
multiplicities, sum \eqref{A:F-negative-large} over $A_z\leq V$
and \eqref{A:F-negative-small} over $A_z>V$.
The term $V/(UA)$ is summable with total
$O_\sigma(U^{-1/2}V^{1-\sigma})$ because $\sigma<0$;
the other two terms are
$O_\sigma(U^{1/2}V^{1-\sigma})$ because $\sigma>-1$.
Thus precisely the same estimate as \eqref{A:H-uniform} holds:
\begin{equation}\label{A:H-negative}
 |H_\ell(\sigma+i\tau)|
       \ll_\sigma U^{1/2}(U+|\kappa_\ell|)^{1-\sigma}.
\end{equation}
The series defining $H_\ell$ converges normally in this strip,
since its summands at large $A_z$ have two additional inverse
powers of $A_z$.

At $s=0$ the exact hypergeometric identity is
\[
 F_{0,\kappa}(A)=\cos\bigl(\kappa\operatorname{arsinh}(1/A)\bigr).
\]
It follows either by solving the differential equation in the proof
of \eqref{A:F-spherical-integral} at $s=0$, or by its power
series. In particular,
\begin{equation}\label{A:H-at-zero}
 H_0(0)=0,\qquad |H_\ell(0)|\ll1+|\ell|.
\end{equation}
For the second assertion sum
$\min(2,\kappa_\ell^2/A_z^2)$ using the same linear counting
bound. No estimate for the distribution of individual ideal angles
is being assumed in this step.

\emph{2. A uniform spectral expansion with square-root-size remainder.}
Fix a large base height $Y_0$, and take
$T=H=\lfloor\sqrt{Y_0}\rfloor$, fixed as $Y$ varies over
$[Y_0,2Y_0]$. Let $w=w_T$ be either smooth radial cutoff of
Statement A.4.4. Its Mellin transform continues meromorphically by
\[
 \widehat w(s)=-\frac1s\int_0^\infty w'(u)u^s\,du.
\]
It has a simple pole at zero of residue one and satisfies
\eqref{A:w-Mellin-bound} on any fixed nonzero real-part line.
Choose $\sigma=-1/2+\eta$, where $0<\eta<1/2$ is fixed.

The contour estimate \eqref{A:contour-uniform}, with $T=H$,
remains valid on this line. Here are the precise changes to its
proof. The gamma bound \eqref{A:gamma-full-bound} is unchanged,
with $a=1-\eta/2<1$ and $b=3/4-\eta/2$; its numerator gamma
functions still have fixed real parts strictly between $-1/2$
and $-1/4$, and so have no pole on the line. The low-angular and
comparable-angular integrations use only $a<1$ and remain
identical. In the high-angular range, $K\ll H=T$ and the sum is
\[
 K^{1/2-\sigma}\sum_{J\ll K}J^{\sigma+1/2}
       \ll_\eta K,
\]
because $\sigma+1/2=\eta>0$. The exponentially small tails
remain $O(\log(2T))$. These facts give
$O_\eta(T\log(2H))$ for the remaining contour before multiplication
by $Y^\sigma$. All constants are uniform in $Y\in[Y_0,2Y_0]$.
Allowing negative angular indices changes only a factor of two.

We also include angular index zero. Its nonconstant spectral
part obeys the same estimates, using the low-angular range.
The constant eigenfunction contributes
\[
 p'_0p_{0,0}
 \frac{\Gamma((s+1)/2)\Gamma((s-1)/2)}{\Gamma(s/2)^2}.
\]
On the negative contour this is $O_\eta(1)$ in $\tau$, and
its only pole in the traversed strip is at $s=1$.
Its residue in $M_0$ is $4c/3$, where
$c=9L_0^2/(2\pi^2)$. This follows also directly from
\eqref{A:Mellin-geometry} and the already proved
$\sum_{A_z\leq Y}\rho(m_z)=(c/3)Y+O(Y^{2/3})$;
$H_0$ is analytic at one.

There is no contribution at zero from $\widehat w(s)M_\ell(s)$:
the normally convergent gamma series shows $M_\ell(0)=0$
for every $\ell$. For $\ell\ne0$, the factor $1/\Gamma(s/2)$
gives this zero; for $\ell=0$ the nonconstant terms have three
such factors, and the displayed constant term has two.
In the geometric correction the pole of $\widehat w$ does give
$H_\ell(0)$, whose total against coefficients
$|b_\ell|\ll(1+|\ell|)^{-1}$ is $O(H)$ by
\eqref{A:H-at-zero}. The two exceptional double cosets contribute
only $O(\log(2H))$.
Equation \eqref{A:H-negative} bounds the negative-contour
correction by
\[
 O_\eta\left(Y_0^\sigma
       T^{1/2}(T+H)^{1-\sigma}\log(2H)\right)
 =O_\eta\left(Y_0^{1/2+\eta/2}\log Y_0\right).
\]
The remaining $M$ contour costs only
$O_\eta(Y_0^\sigma T\log(2H))$.
The finite-height contour justification from Statement A.4.4
continues without change; the next spectral poles have real part
$-3/2$, outside the traversed strip.

Consequently, for any real trigonometric polynomial
$B(\theta)=\sum_{|\ell|\leq H}b_\ell e^{2\pi i\ell\theta}$
with $|b_\ell|\ll(1+|\ell|)^{-1}$, the weighted count
\[
 \mathcal N_{w,B}(Y)=\sum_z\sum_{N\mathfrak a=m_z}
                  w(A_z/Y)B(\theta(\mathfrak a))
\]
has the expansion, uniformly on $[Y_0,2Y_0]$,
\begin{equation}\label{A:sqrt-remainder-expansion}
 \mathcal N_{w,B}(Y)=\frac{b_0c}{3}\widehat w(1)Y
 +\frac14\sum_{j\geq1}p'_j
  \left(r_j\widehat w(s_j)Q_{j,H}Y^{s_j}
       +\overline{r_j\widehat w(s_j)Q_{j,H}}Y^{1-s_j}\right)
 +O_\eta\left(Y_0^{1/2+\eta/2}\log Y_0\right).
\end{equation}
Here $Q_{j,H}=\sum b_\ell D_{s_j}(0)
p_{j,\ell}/D_{s_j}(\kappa_\ell)$ uses the unnormalized source
periods. The conjugation in the second term is valid because
the eigenfunctions are real and $b_{-\ell}=\overline{b_\ell}$.
The coefficients, including $\widehat w(s_j)$ and $Q_{j,H}$,
are fixed throughout this $Y$ interval.

\emph{3. The large sieve and its hypotheses.}
Theorem 3 of Lekkas--Voskou \cite{LekkasVoskou2024}, with weight
$m=0$, states that for a fixed cofinite Fuchsian group and a
fixed closed geodesic, arbitrary complex coefficients $d_j$, and
$\delta$-separated points $y_\nu\in[Y_0,2Y_0]$,
\[
 \sum_\nu\left|\sum_{|t_j|\leq V}d_jy_\nu^{it_j}
                    \int_lu_j\right|^2
 \ll (V+Y_0\delta^{-1})\sum_{|t_j|\leq V}|d_j|^2.
\]
The continuous-spectrum coefficients vanish on a compact quotient.
The same statement is Theorem 1.3 of \cite{LekkasPetridis2026}.
Apply it on the fixed compact torsion-free cover from Statement A.3.1,
with zero coefficients for all eigenfunctions not lifted from
$\Gamma^+$. The covering and period normalization factors are
fixed and were verified above. Apply it separately to the two
target geodesics and use Cauchy--Schwarz for their sum $p_j'$.
There is no requirement that the target equal the source geodesic
or that the coefficients $d_j$ be periods.

For a finite spectral sum, take equally spaced $y_\nu$ with
mesh $h$, choose $\delta=h/2$, multiply the inequality by $h$,
and let $h\to0$. The Riemann sums converge, giving
\begin{equation}\label{A:continuous-large-sieve}
 \int_{Y_0}^{2Y_0}
       \left|\sum_jd_jY^{it_j}p_j'\right|^2dY
       \ll Y_0\sum_j|d_j|^2.
\end{equation}
For the band $J\leq1+t_j<2J$, take
$d_j=r_j\widehat w(s_j)Q_{j,H}$.
By Statement A.4.3, Stirling's formula, and
\eqref{A:w-Mellin-bound},
\[
 \sum_{j\text{ in band}}|d_j|^2
 \ll J^{-1}(1+J/T)^{-2N}\sum_{j\text{ in band}}|Q_{j,H}|^2
 \ll\log^2(2H)(1+J/T)^{-2N}.
\]
The factor $Y^{1/2}$ in \eqref{A:sqrt-remainder-expansion}
therefore gives $L^2([Y_0,2Y_0])$ norm at most
$O(Y_0\log(2H)(1+J/T)^{-N})$ for this band.
Minkowski's inequality over the dyadic bands, and then passage
to the absolutely convergent spectral series, prove
\begin{equation}\label{A:smoothed-mean-square}
 \int_{Y_0}^{2Y_0}
 \left|\mathcal N_{w,B}(Y)-\frac{b_0c}{3}\widehat w(1)Y\right|^2dY
 \ll_\eta Y_0^2\log^4Y_0+Y_0^{2+\eta}\log^2Y_0.
\end{equation}

\emph{4. Sharp arc and sharp height.}
For the fixed arc $I=[0,\lambda_0]$ take its Selberg majorant
and minorant trigonometric polynomials $B_H^+,B_H^-$.
The exact elementary approximation used here is
\cite[equation (4.11)]{LesterMatomakiRadziwill2015}:
for every interval in $\mathbb R/\mathbb Z$ and integer $H\geq1$,
\[
 B_H^-\leq\mathbf1_I\leq B_H^+,\qquad
 b_0^\pm=|I|\pm\frac1{H+1},\qquad
 |b_\ell^\pm|\ll(1+|\ell|)^{-1}\quad(0<|\ell|\leq H).
\]
Endpoint half-value conventions can change the sum only by the
already proved $O(\log Y_0)$ Pell contribution.
These polynomials are real, have degree at most $H$, and are
fixed throughout the dyadic $Y$ interval. In particular $B_H^+\geq0$.
Use the radial cutoffs $w^\pm$ of Statement A.4.4 with the same
fixed $T=\lfloor\sqrt{Y_0}\rfloor$. Then
\[
 \mathcal N_{w^-,B_H^-}(Y)
 \leq\mathcal A(Y)
 \leq\mathcal N_{w^+,B_H^+}(Y)+O(\log Y_0),
\]
where $\mathcal A(Y)$ is the exact sharp ideal-arc count with
$A_z\leq Y$. The lower inequality is valid even where
$B_H^-$ is negative: if $w^-\ne0$ the radial sharp indicator
is one, and multiplication by $0\leq w^-\leq1$ still leaves
$w^-B_H^-\leq\mathbf1_I$. Outside that radial range it is zero.
The upper inequality follows from nonnegativity of $B_H^+$.

Both main terms differ from $\lambda_0cY/3$ by
$O(Y_0/H+Y_0/T)=O(\sqrt{Y_0})$.
Thus \eqref{A:smoothed-mean-square} for the two fixed pairs
$(w^-,B_H^-)$ and $(w^+,B_H^+)$ proves
\[
 \int_{Y_0}^{2Y_0}|\mathcal A(Y)-\lambda_0cY/3|^2dY
       \ll_\epsilon Y_0^{2+\epsilon},
\]
after choosing $\eta$ smaller than $\epsilon$ and absorbing
logarithms. Finally $Y=(9x^2-1)^{1/2}$ has derivative bounded
above and below by positive constants for $x\geq1$;
two dyadic $Y$ intervals cover the image of $[X,2X]$.
The difference between $\lambda_0cY/3$ and $C_0x$ is
$O(x^{-1})$, and Statement A.1.3 adds $O(\log(2x))$.
These terms have mean square $O(X\log^2X)$, proving
\eqref{A:sharp-mean-square}.
\end{proof}
